\subsection{有限幂零群}

定理 4. 设 G 为有限群。下列条件等价：

\begin{enumerate}
\def\labelenumi{(\alph{enumi})}
\item
  G 是幂零的。
\item
  G 是一些 p-群的乘积。
\item
  对每个素数 \(p\) ，存在一个正规的西罗 \(p\) -子群（ \(G\) 的）。
\item
  (b) \(\Rightarrow\) (a)（第 5 号，定理 1 的推论）。
\end{enumerate}

假设 (a) 成立，设 P 为 G 的一个西罗 p-子群。若 N 为 P 在 G 中的正规化子，则定理 3 的推论 1 表明 N 是其自身的正规化子。由 §6 第 3 号命题 8 的推论，这表明 N = G。因此 (a) ⇒ (c)。

假设 (c) 成立，设 I 为整除 \(\operatorname{Card}(G)\) 的素数所成的集合。对所有 \(p \in I\) ，设 \(P_p\) 是一个正规西罗 \(p\) -子群（ \(G\) 的）。对所有 \(p \neq q\) ， \(P_p \cap P_q\) 缩减为 \(e\) ，因为它既是一个 \(p\) -群又是一个 \(q\) -群，因此 \(P_p\) 与 \(P_q\) 彼此中心化（§4，第 9 号，命题 15）。设 \(\phi\) 是 \(\prod_{p \in I} P_p\) 到 \(G\) 的典范同态（§4，第 9 号，命题 12）。同态 \(\phi\) 是满射的，这由第 6 号的注记得出。由于 \(\operatorname{Card}\left(\prod_{p \in I} P_p\right) = \operatorname{Card}(G)\) ，由此可得 \(\phi\) 是双射。

注记。(1) 设 G 为有限群，p 为素数。由第 6 号定理 3 (a) 和第 6 号定理 2，下列条件等价：

\begin{enumerate}
\def\labelenumi{(\roman{enumi})}
\item
  存在一个正规的西罗 \(p\) -子群（ \(\mathbf{G}\) 的）；
\item
  每个西罗 \(p\) -子群（ \(\mathbf{G}\) 的）都是正规的；
\item
  G 仅有一个西罗 p-子群。
\end{enumerate}

\begin{enumerate}
\def\labelenumi{(\arabic{enumi})}
\setcounter{enumi}{1}
\item
  设 \(\mathbf{G}\) 是一个有限幂零群。设 \(\mathbf{I}\) 为 \(\operatorname{Card}(\mathbf{G})\) 的素因子所成的集合。由定理 4 和注记 1， \(\mathbf{G} = \prod_{p \in \mathbf{I}} \mathbf{G}_p\) ，其中 \(\mathbf{G}_p\) 是唯一的西罗 \(p\) -群（ \(\mathbf{G}\) 的）。
\item
  将定理 4 应用于交换群，就得到交换有限群分解为准素分支之积的分解，这将在第七章中从另一个角度加以研究。
\end{enumerate}

例子。群 \(\mathfrak{S}_3\) 的阶为 6。它包含一个 3 阶的正规西罗 3-子群：群 \(\mathfrak{A}_3\) 。它包含三个 2 阶的西罗 2-子群：群 \(\{e, \tau\}\) ，其中 \(\tau\) 是一个对换。因此群 \(\mathfrak{S}_3\) 不是幂零的。

